% 図: ネットワーク層の整理図（OSI と TCP/IP の対応）
% 用途: % 図: ネットワーク層の整理図（OSI と TCP/IP の対応）
% 用途: % 図: ネットワーク層の整理図（OSI と TCP/IP の対応）
% 用途: % 図: ネットワーク層の整理図（OSI と TCP/IP の対応）
% 用途: \input{assets/diagrams/network_layers}
\begin{tikzpicture}[node distance=1mm and 4mm, every node/.style={font=\footnotesize}]

  % --- OSI 7層 ---
  \node[keynode, minimum width=32mm] (osi7) {7. アプリケーション};
  \node[keynode, minimum width=32mm, below=of osi7] (osi6) {6. プレゼンテーション};
  \node[keynode, minimum width=32mm, below=of osi6] (osi5) {5. セション};
  \node[keynode, minimum width=32mm, below=of osi5] (osi4) {4. トランスポート};
  \node[keynode, minimum width=32mm, below=of osi4] (osi3) {3. ネットワーク};
  \node[keynode, minimum width=32mm, below=of osi3] (osi2) {2. データリンク};
  \node[keynode, minimum width=32mm, below=of osi2] (osi1) {1. 物理};

  % --- TCP/IP 4層（右側にまとめて対応） ---
  \node[boxnode, minimum width=36mm, right=10mm of osi6] (ta) {アプリケーション層\\\scriptsize HTTP / SMTP / DNS};
  \node[boxnode, minimum width=36mm, below=3mm of ta]   (tt) {トランスポート層\\\scriptsize TCP / UDP};
  \node[boxnode, minimum width=36mm, below=3mm of tt]   (tn) {インターネット層\\\scriptsize IP / ICMP};
  \node[boxnode, minimum width=36mm, below=3mm of tn]   (tl) {ネットワーク\\インタフェース層\\\scriptsize Ethernet / Wi-Fi};

  % --- 対応線 ---
  \draw[arr2] (osi7.east) -- (ta.west);
  \draw[arr2] (osi6.east) -- (ta.west);
  \draw[arr2] (osi5.east) -- (ta.west);
  \draw[arr2] (osi4.east) -- (tt.west);
  \draw[arr2] (osi3.east) -- (tn.west);
  \draw[arr2] (osi2.east) -- (tl.west);
  \draw[arr2] (osi1.east) -- (tl.west);

  \node[subnode, below=6mm of osi1, minimum width=88mm, text width=82mm, align=center]
    {左が理論の整理（OSI）／右が実運用の整理（TCP/IP）。番号が上ほど「人に近い」。};

\end{tikzpicture}

\begin{tikzpicture}[node distance=1mm and 4mm, every node/.style={font=\footnotesize}]

  % --- OSI 7層 ---
  \node[keynode, minimum width=32mm] (osi7) {7. アプリケーション};
  \node[keynode, minimum width=32mm, below=of osi7] (osi6) {6. プレゼンテーション};
  \node[keynode, minimum width=32mm, below=of osi6] (osi5) {5. セション};
  \node[keynode, minimum width=32mm, below=of osi5] (osi4) {4. トランスポート};
  \node[keynode, minimum width=32mm, below=of osi4] (osi3) {3. ネットワーク};
  \node[keynode, minimum width=32mm, below=of osi3] (osi2) {2. データリンク};
  \node[keynode, minimum width=32mm, below=of osi2] (osi1) {1. 物理};

  % --- TCP/IP 4層（右側にまとめて対応） ---
  \node[boxnode, minimum width=36mm, right=10mm of osi6] (ta) {アプリケーション層\\\scriptsize HTTP / SMTP / DNS};
  \node[boxnode, minimum width=36mm, below=3mm of ta]   (tt) {トランスポート層\\\scriptsize TCP / UDP};
  \node[boxnode, minimum width=36mm, below=3mm of tt]   (tn) {インターネット層\\\scriptsize IP / ICMP};
  \node[boxnode, minimum width=36mm, below=3mm of tn]   (tl) {ネットワーク\\インタフェース層\\\scriptsize Ethernet / Wi-Fi};

  % --- 対応線 ---
  \draw[arr2] (osi7.east) -- (ta.west);
  \draw[arr2] (osi6.east) -- (ta.west);
  \draw[arr2] (osi5.east) -- (ta.west);
  \draw[arr2] (osi4.east) -- (tt.west);
  \draw[arr2] (osi3.east) -- (tn.west);
  \draw[arr2] (osi2.east) -- (tl.west);
  \draw[arr2] (osi1.east) -- (tl.west);

  \node[subnode, below=6mm of osi1, minimum width=88mm, text width=82mm, align=center]
    {左が理論の整理（OSI）／右が実運用の整理（TCP/IP）。番号が上ほど「人に近い」。};

\end{tikzpicture}

\begin{tikzpicture}[node distance=1mm and 4mm, every node/.style={font=\footnotesize}]

  % --- OSI 7層 ---
  \node[keynode, minimum width=32mm] (osi7) {7. アプリケーション};
  \node[keynode, minimum width=32mm, below=of osi7] (osi6) {6. プレゼンテーション};
  \node[keynode, minimum width=32mm, below=of osi6] (osi5) {5. セション};
  \node[keynode, minimum width=32mm, below=of osi5] (osi4) {4. トランスポート};
  \node[keynode, minimum width=32mm, below=of osi4] (osi3) {3. ネットワーク};
  \node[keynode, minimum width=32mm, below=of osi3] (osi2) {2. データリンク};
  \node[keynode, minimum width=32mm, below=of osi2] (osi1) {1. 物理};

  % --- TCP/IP 4層（右側にまとめて対応） ---
  \node[boxnode, minimum width=36mm, right=10mm of osi6] (ta) {アプリケーション層\\\scriptsize HTTP / SMTP / DNS};
  \node[boxnode, minimum width=36mm, below=3mm of ta]   (tt) {トランスポート層\\\scriptsize TCP / UDP};
  \node[boxnode, minimum width=36mm, below=3mm of tt]   (tn) {インターネット層\\\scriptsize IP / ICMP};
  \node[boxnode, minimum width=36mm, below=3mm of tn]   (tl) {ネットワーク\\インタフェース層\\\scriptsize Ethernet / Wi-Fi};

  % --- 対応線 ---
  \draw[arr2] (osi7.east) -- (ta.west);
  \draw[arr2] (osi6.east) -- (ta.west);
  \draw[arr2] (osi5.east) -- (ta.west);
  \draw[arr2] (osi4.east) -- (tt.west);
  \draw[arr2] (osi3.east) -- (tn.west);
  \draw[arr2] (osi2.east) -- (tl.west);
  \draw[arr2] (osi1.east) -- (tl.west);

  \node[subnode, below=6mm of osi1, minimum width=88mm, text width=82mm, align=center]
    {左が理論の整理（OSI）／右が実運用の整理（TCP/IP）。番号が上ほど「人に近い」。};

\end{tikzpicture}

\begin{tikzpicture}[node distance=1mm and 4mm, every node/.style={font=\footnotesize}]

  % --- OSI 7層 ---
  \node[keynode, minimum width=32mm] (osi7) {7. アプリケーション};
  \node[keynode, minimum width=32mm, below=of osi7] (osi6) {6. プレゼンテーション};
  \node[keynode, minimum width=32mm, below=of osi6] (osi5) {5. セション};
  \node[keynode, minimum width=32mm, below=of osi5] (osi4) {4. トランスポート};
  \node[keynode, minimum width=32mm, below=of osi4] (osi3) {3. ネットワーク};
  \node[keynode, minimum width=32mm, below=of osi3] (osi2) {2. データリンク};
  \node[keynode, minimum width=32mm, below=of osi2] (osi1) {1. 物理};

  % --- TCP/IP 4層（右側にまとめて対応） ---
  \node[boxnode, minimum width=36mm, right=10mm of osi6] (ta) {アプリケーション層\\\scriptsize HTTP / SMTP / DNS};
  \node[boxnode, minimum width=36mm, below=3mm of ta]   (tt) {トランスポート層\\\scriptsize TCP / UDP};
  \node[boxnode, minimum width=36mm, below=3mm of tt]   (tn) {インターネット層\\\scriptsize IP / ICMP};
  \node[boxnode, minimum width=36mm, below=3mm of tn]   (tl) {ネットワーク\\インタフェース層\\\scriptsize Ethernet / Wi-Fi};

  % --- 対応線 ---
  \draw[arr2] (osi7.east) -- (ta.west);
  \draw[arr2] (osi6.east) -- (ta.west);
  \draw[arr2] (osi5.east) -- (ta.west);
  \draw[arr2] (osi4.east) -- (tt.west);
  \draw[arr2] (osi3.east) -- (tn.west);
  \draw[arr2] (osi2.east) -- (tl.west);
  \draw[arr2] (osi1.east) -- (tl.west);

  \node[subnode, below=6mm of osi1, minimum width=88mm, text width=82mm, align=center]
    {左が理論の整理（OSI）／右が実運用の整理（TCP/IP）。番号が上ほど「人に近い」。};

\end{tikzpicture}
